% ===========================================================================
\chapter{Asking questions}\label{ch:queries}

Once the database is at a fixpoint, a question is easy. \code{?- sees(ann,
Who).} asks for every value of \code{Who} that makes the atom a fact, and the
facts are all there: match the question against each candidate fact, and
every match is an answer.

\begin{newfile}{query.nex}
-- query.nex: answering questions.

intern Database

-- Every way of binding the variables of `q` so that it is a fact in `db`.
function answers(q: Atom, db: Database): Array[Bindings]
do
  result := []
  let none := create Bindings.empty
  across db.candidates(q, none) as f do
    let b := match_atom(q, f, none)
    if b /= nil then
      result.add(b)
    end
  end
ensure
  all_true: all_facts(q, result, db)
end
\end{newfile}

The postcondition says what an answer is: the question, with the answer's
bindings substituted, is a fact in the database.

\begin{continued}{query.nex}
-- Is `q`, under each of `found`, a fact in `db`?
function all_facts(q: Atom, found: Array[Bindings], db: Database): Boolean
do
  result := true
  across found as b do
    if not db.contains(substitute(q, b)) then
      result := false
    end
  end
end
\end{continued}

For printing, each answer becomes its bindings as text, sorted, so that the
same question always gives the same output. A question without variables,
\code{?- friends(ann, bo).}, has a yes-or-no answer.

\begin{continued}{query.nex}
-- The answers to `q` as lines of text, in order: "yes" or "no" for a
-- question without variables.
function answer_lines(q: Atom, db: Database): Array[String]
do
  let found := answers(q, db)
  if q.is_ground then
    result := [when found.is_empty then "no" else "yes" end]
  elseif found.is_empty then
    result := ["no answers"]
  else
    result := found.map(fn (b: Bindings): String do
      result := b.text
    end).sort
  end
end
\end{continued}

\section{Trying it out}

\begin{tryit}{try_queries.nex}
intern Parser
intern Evaluate
intern Query
intern io/Path

match parse_program(create Path.make(program.working_directory + "/social.dl").read_text) of
  Ok as ok then
    do
      let db := create Database.make
      across ok.value.facts as f do
        db.add(f)
      end
      (create Naive_Evaluator.make).run(ok.value.rules, db)
      across ok.value.queries as q do
        console.print_line("?- " + q.text + ".")
        across answer_lines(q, db) as line do
          console.print_line("   " + line)
        end
      end
      match parse_program("?- friends(ann, bo). ?- friends(ann, cy). ?- sees(eli, X).") of
        Ok as more then
          across more.value.queries as q do
            console.print_line("?- " + q.text + ".   " + answer_lines(q, db).to_string)
          end
        Err as e then console.print_line(e.error)
      end
    end
  Err as e then console.print_line(e.error)
end
\end{tryit}
\tryoutput{try_queries}

Ann is among the people whose posts reach her. She follows Bo, who follows
her back, so her own posts come back to her through his reshares. The rules
say nothing to rule that out, and the engine reports exactly what they say.
